\section{Experimental Setup}
\label{sec:setup}

\subsection{Models}

\paragraph{LLMs.}
We evaluate ten frontier LLMs spanning six families: three Gemini~3 models via the Gemini API (Gemini~3.1~Pro, Gemini~3~Flash, Gemini~3.1~Flash-Lite) and seven open-weight models served via OpenRouter (DeepSeek-R1, DeepSeek-V4-Flash, Qwen3.6-27B, Qwen3.6-35B-A3B, GLM-4.7, Kimi-K2.6, and MiniMax-M2.7).
Together they span dense and mixture-of-experts architectures, reasoning and instruct variants, and a wide capability and cost range.
Gemini~3.1~Flash-Lite provides a non-reasoning baseline; the other models support extended reasoning (chain-of-thought tokens \citep{wei2022cot}, billed at the output token rate).

\paragraph{Embedding models.}
We evaluate 26 text embedding models from 118M to 14B parameters. The multi-model families are multilingual-E5 \citep{wang2024textembeddingsweaklysupervisedcontrastive,wang2024improvingtextembeddingslarge}, Qwen3-Embedding \citep{qwen2025qwen3embedding}, F2LLM-v2 \citep{zhang2026f2llmv2}, Jina-v5 \citep{akram2026jinaembeddingsv5}, and GTE-Qwen2 \citep{li2023generaltextembeddingsmultistage}. We also include models from Tencent \citep{hu2025kalmv2}, NVIDIA \citep{babakhin2025llamaembednemotron}, Salesforce \citep{meng2024sfremb2}, Snowflake \citep{yu2024arctic}, and Google \citep{vera2025embeddinggemma}, together with E5-Mistral \citep{wang2024improvingtextembeddingslarge}, GritLM \citep{muennighoff2025generativerepresentationalinstructiontuning}, Linq-Embed-Mistral \citep{choi2024linq}, BGE-M3 \citep{chen2025m3embeddingmultilingualitymultifunctionalitymultigranularity}, and Octen-8B.
All embedding models run locally on a single NVIDIA H100 80GB HBM3 GPU.
Table~\ref{tab:main_results} lists every model with its score and cost; full details appear in Appendix~\ref{app:models_detail}.


\subsection{Tasks and Evaluation Protocol}
\label{sec:tasks}

We introduce \textbf{\mteblm}, a 37-task benchmark covering classification, STS, clustering, pair classification, and retrieval, implemented with the MTEB framework \citep{muennighoff2023mteb}.
It follows MTEB's task and result interfaces, allowing new models and datasets to use the same evaluation pipeline.
Each \mteblm{} task is a new LLM-specific evaluation dataset derived from the corresponding original MTEB task (fixed seed\,=\,42; Appendix~\ref{app:datasets}), ensuring all submissions are evaluated on identical data.

\paragraph{Why these tasks.}
We use subsets rather than MTEB(eng, v2) because generative evaluation is orders of magnitude more expensive: one pass over the full suite would cost hundreds of dollars per LLM, and corpus-in-context retrieval requires the corpus to fit in the prompt.
Tasks are selected to span all five categories, cover multiple languages and domains, and keep retrieval corpora small enough for shorter-context models.
MTEB's task-selection analysis was also designed around embedding models, so v2 is not automatically the right frame for a cross-paradigm comparison.

We compare complete deployment pipelines: embedding models use kNN for classification, cosine similarity for STS and retrieval, and $k$-means for clustering; LLMs receive zero-shot prompts.
This comparison reflects a common deployment setting in which practitioners consider LLMs because labelled data are scarce.
Our few-shot ablation (\S\ref{sec:fewshot}) tests whether a small number of in-context examples narrows the resulting gap.

\paragraph{Classification (8 tasks).}
Sentiment (IMDB, ToxicConversations, TweetSentiment), intent detection (Banking77: 77 classes; MassiveIntent: 60 intents), and multilingual classification (AmazonCounterfactual, MTOPDomain, MassiveScenario).
\emph{LLM:} Zero-shot structured output for each test split sample.
\emph{Embedding:} kNN trained on embeddings on the labelled training split, then used for test split predictions.
\emph{Metric:} Accuracy.

\paragraph{Semantic Textual Similarity (10 tasks).}
English benchmarks (STSBenchmark, SICK-R, STS12--16), a biomedical benchmark (BIOSSES), and multilingual benchmarks (STS17, STS22v2).
\emph{LLM:} Rates similarity on each task's native scale, returning a float.
\emph{Embedding:} Cosine similarity between embeddings.
\emph{Metric:} Spearman correlation.

\paragraph{Clustering (9 tasks).}
Scientific papers (ArXiv, BioRxiv, MedRxiv) and online discussions (Reddit, StackExchange, TwentyNewsgroups).
\emph{LLM:} All documents in a single prompt with sequential identifiers; the model receives the ground-truth cluster count $k$ and returns cluster assignments as a JSON array.
\emph{Embedding:} $k$-means on document embeddings.
\emph{Metric:} V-measure.

\paragraph{Pair Classification (4 tasks).}
Duplicate detection (SprintDuplicateQuestions), paraphrase identification (TwitterURLCorpus), legal matching (LegalBenchPC), and entailment (RTE3, multilingual).
\emph{LLM:} Binary output per pair with reasoning.
\emph{Embedding:} Cosine similarity with threshold sweep (MTEB convention; best threshold chosen post-hoc on the test set).
\emph{Metric:} Average precision and accuracy.

\paragraph{Retrieval (6 tasks).}
A domain-diverse suite: legal (AILAStatutes; LegalBench consumer-contracts QA \citep{guha2023legalbench}), finance (HC3-Finance \citep{guo2023hc3}), public-health QA, French QA (FQuAD \citep{dhoffschmidt2020fquad}), and Danish social media (TwitterHjerne \citep{holm2024danoliterate}).
\emph{LLM:} Corpus-in-context \citep{lee2024longcontextlanguagemodelssubsume}: the full corpus is placed in the prompt with sequential identifiers; prompt caching \citep{gim2024promptcache} amortises the corpus prefix across queries.
Corpora are deliberately small (82--415 documents) so that models with shorter context windows can be evaluated on identical data; this keeps the whole corpus readable in one prompt, which production-scale retrieval does not permit (see \S\ref{sec:limitations}).
\emph{Embedding:} Cosine similarity ranking.
\emph{Metric:} Recall@1 (nDCG@$k$ results in released data files).


\subsection{Cost Methodology}
\label{sec:cost}

\paragraph{Embedding costs.}
We measure maximum-throughput token processing on an H100 80GB at sequence length~512 (largest fitting batch) and compute cost as $r_{\text{GPU}} / T$, where $r_{\text{GPU}} = \$2.49$/hr (Lambda Labs H100 spot, March 2026) and $T$ is tokens per hour.
Costs range from \$0.001 (mE5-small) to \$0.22 (F2LLM-v2-14B); full per-model throughput and cost data appear in Table~\ref{tab:embedding_throughput} (Appendix~\ref{app:emb_throughput}).

\paragraph{LLM costs.}
Token usage is extracted from each provider's \texttt{usage\_stats}. We bill every non-input (generated) token at the output rate:
\begin{equation}
\text{Cost} = (\text{input} - \text{cached}) \cdot r_{\text{in}} + \text{cached} \cdot r_{\text{cache}} + (\text{total} - \text{input}) \cdot r_{\text{out}}
\label{eq:cost}
\end{equation}
where $r_{\text{cache}} = r_{\text{in}} / 10$. Reasoning (``thinking'') tokens are part of the generated total and are billed at the output rate, matching provider billing; providers report them differently, and we attribute each model's reasoning share from its usage statistics (Appendix~\ref{app:cost_detail}).
Gemini models are priced at Gemini API rates and open-weight models at OpenRouter public rates (March/June 2026).
Costs range from \$3.16 (DeepSeek-V4-Flash) to \$154.14 (Gemini~3.1~Pro); the complete token-level breakdown is in Table~\ref{tab:llm_tokens} (Appendix~\ref{app:tokens_table}).


\subsection{Throughput Methodology}
\label{sec:throughput}

We serve both paradigms on the same GPU, giving the throughput comparison common hardware and removing API rate limits.
This differs from our cost accounting (\S\ref{sec:cost}), which uses API rates for LLMs and GPU-rental rates for embeddings.
We serve the two open-weight LLMs that fit on a single H100 (Qwen3.6-27B dense, Qwen3.6-35B-A3B MoE) under vLLM, and run embedding models batched on the same GPU, measuring both in tokens per second (configuration in Appendix~\ref{app:cost_detail}).
Gemini is API-only, and larger open models such as DeepSeek-R1 need multiple GPUs.
This end-to-end comparison uses each paradigm's standard inference procedure: encoder batching for embeddings and autoregressive decoding for LLMs (Appendix~\ref{app:extended_limitations}).
Per-model embedding throughput is reported in Table~\ref{tab:embedding_throughput} (Appendix~\ref{app:emb_throughput}); the comparison appears in Figure~\ref{fig:throughput} (Appendix~\ref{app:figures}).
